\chapter{Taux et change}\label{ch:iv2:rates-and-fx}

\begin{ivtrains}
Courbes et ce qu'on peut en tirer, couverture au DV01, swaps et contrats à terme, modèles de taux court et swaptions, puis forwards de change et options de change. Sauf indication contraire dans un énoncé, les rendements sont composés annuellement, et chaque niveau présenté ici est inventé pour les besoins du calcul.
\end{ivtrains}

\section{Courbes à toute vitesse}

\begin{question}[$\star$ \iqroles{QR, QT, Trader, Risque} \iqfirm{banque, hedge fund} \iqlevel{junior}]\label{iq:iv2:rates-and-fx:1}
Une obligation zéro-coupon à cinq ans rapporte 4~\%, avec une composition annuelle. Quel est son prix pour 100~\$~? Quel est le prix si les 4~\% sont composés en continu~?
\end{question}

\begin{question}[$\star$ \iqroles{QR, QT, Trader, Risque} \iqfirm{banque, hedge fund} \iqlevel{junior}]\label{iq:iv2:rates-and-fx:2}
Une obligation à un an verse un coupon de 5~\$ et 100~\$ à l'échéance, et son rendement est de 4~\%. Quel est son prix~? Sans calcul, dites quel serait son prix à un rendement de 5~\%, et pourquoi.
\end{question}

\begin{question}[$\star$ \iqroles{QR, QT, Trader, Risque} \iqfirm{banque, hedge fund} \iqlevel{junior}]\label{iq:iv2:rates-and-fx:3}
Un zéro-coupon à cinq ans rapporte 4~\%. Quelles sont ses durations de Macaulay et modifiée~? De combien varie le prix de 100~\$ si les rendements montent de 10 points de base~?
\end{question}

\begin{question}[$\star$ \iqroles{QR, QT, Trader, Risque} \iqfirm{banque, hedge fund} \iqlevel{junior}]\label{iq:iv2:rates-and-fx:4}
Vous détenez 10 millions de dollars d'une obligation cotée au pair, de duration modifiée 8. Quel est votre DV01~? Combien gagnez-vous si les rendements baissent de 25 points de base, en négligeant la convexité~?
\end{question}

\begin{question}[$\star$ \iqroles{QR, QT, Trader, Risque} \iqfirm{banque, hedge fund} \iqlevel{junior}]\label{iq:iv2:rates-and-fx:5}
Les zéros à un an rapportent 3~\% et les zéros à deux ans 4~\%. Quel est le taux à un an dans un an, sous forme de forward~? Que fait un emprunteur pour le fixer dès aujourd'hui~?
\end{question}

\begin{question}[$\star$ \iqroles{QT, Trader, Risque} \iqfirm{banque, hedge fund} \iqlevel{junior}]\label{iq:iv2:rates-and-fx:6}
Vous êtes acheteur de 10 millions de dollars d'une obligation à cinq ans de duration modifiée 4.5, au pair. Vous couvrez avec une obligation à dix ans de duration modifiée 8, également au pair. Quel montant vendez-vous à découvert~? À quoi restez-vous exposé~?
\end{question}

\begin{question}[$\star$ \iqroles{QR, QT, Trader, Risque} \iqfirm{banque, hedge fund} \iqlevel{junior}]\label{iq:iv2:rates-and-fx:7}
Vous payez le taux fixe de 4~\% sur un swap de 10 millions de dollars à cinq ans, et vous recevez le taux variable. Il a un DV01 d'environ 4\,500~\$. Gagnez-vous ou perdez-vous si les taux montent~? À quelle position obligataire cela ressemble-t-il~? Que gagnez-vous si le taux du swap monte de 10 points de base~?
\end{question}

\begin{question}[$\star$ \iqroles{QT, Trader, Analyste} \iqfirm{banque, hedge fund} \iqlevel{junior}]\label{iq:iv2:rates-and-fx:8}
Un bon du Trésor à trois mois se négocie avec une décote de 4~\%, sur la base 90/360. Quel est son prix pour 100~\$~? Quel est le rendement monétaire simple sur la même base~? Pourquoi est-il supérieur à 4~\%~?
\end{question}

\begin{question}[$\star$ \iqroles{QT, Trader, QR} \iqfirm{banque, hedge fund} \iqlevel{junior}]\label{iq:iv2:rates-and-fx:9}
L'EURUSD vaut 1.10. Le taux du dollar à un an est de 5~\%, celui de l'euro à un an de 3~\%, tous deux simples. Quel est le forward à un an~? Donnez les points de terme. Le forward annonce-t-il une hausse de l'euro~?
\end{question}

\begin{question}[$\star$ \iqroles{QT, Trader, Analyste} \iqfirm{banque, hedge fund} \iqlevel{junior}]\label{iq:iv2:rates-and-fx:10}
L'EURUSD vaut 1.10 et l'USDJPY vaut 150. Quel est l'EURJPY~? Si l'EURJPY est coté 166, que faites-vous~? Que devez-vous vérifier avant de parler de profit sans risque~?
\end{question}

\begin{question}[$\star$ \iqroles{QT, Trader, QR} \iqfirm{banque, hedge fund} \iqlevel{junior}]\label{iq:iv2:rates-and-fx:11}
Vous détenez un call sur 1 million d'euros contre dollars, à un strike de 1.10. Quelle option cela représente-t-il vu de l'autre côté, et sur combien de dollars~? À quel strike en euros par dollar~?
\end{question}

\begin{question}[$\star$ \iqroles{QT, Trader, QR} \iqfirm{banque, hedge fund} \iqlevel{junior}]\label{iq:iv2:rates-and-fx:12}
Vous empruntez des euros à 3~\% et prêtez des dollars à 5~\%, sans couverture, avec un EURUSD à 1.10. À quel niveau l'EURUSD doit-il se trouver dans un an pour que vous soyez à l'équilibre, en taux simples~? Quel lien cela a-t-il avec le forward~? Pourquoi le nom du trade ne vous dit-il rien de son risque~?
\end{question}

\begin{question}[$\star$ \iqroles{QT, Trader, QR, Risque} \iqfirm{banque, hedge fund} \iqlevel{junior}]\label{iq:iv2:rates-and-fx:13}
Une obligation a une duration modifiée de 8 et une convexité de 80. Estimez la variation de prix en pourcentage pour une hausse et pour une baisse de 100 points de base. Laquelle est la plus grande, et pourquoi les investisseurs apprécient-ils cela~?
\end{question}

\section{Couvertures et trades de courbe}

\begin{question}[$\star\star$ \iqroles{QR, QT, QD} \iqfirm{banque, hedge fund} \iqlevel{junior}]\label{iq:iv2:rates-and-fx:14}
Les rendements au pair annuels à un an et à deux ans valent 3~\% et 4~\%. Bootstrappez les facteurs d'actualisation. Quel est le taux zéro à deux ans~? Quel est le taux forward entre les années un et deux~?
\end{question}

\begin{question}[$\star\star$ \iqroles{QR, QT, QD} \iqfirm{banque, hedge fund} \iqlevel{junior}]\label{iq:iv2:rates-and-fx:15}
Les taux zéro à un, deux et trois ans valent 3~\%, 4~\% et 5~\%. Quel est le taux de swap au pair à trois ans, avec des paiements annuels~? Quel est le swap à deux ans qui débute dans un an~? Pourquoi la jambe variable n'a-t-elle pas besoin de modèle~?
\end{question}

\begin{question}[$\star\star$ \iqroles{QT, Trader, QR} \iqfirm{banque, hedge fund} \iqlevel{junior, confirmé}]\label{iq:iv2:rates-and-fx:16}
Sur une courbe plate à 4~\%, vous comparez une obligation au pair à cinq ans à 4~\% avec un barbell formé d'une obligation au pair à deux ans et d'une obligation au pair à dix ans, toutes deux à 4~\%, de même duration. Comment pondérez-vous les deux~? Quelles sont leurs convexités~? Que rapporte chacune lors d'un mouvement parallèle de 100 points de base à la hausse ou à la baisse~? Alors pourquoi quelqu'un détient-il l'obligation bullet, sachant qu'une vraie courbe n'est pas plate~?
\end{question}

\begin{question}[$\star\star$ \iqroles{QT, Trader} \iqfirm{banque, hedge fund} \iqlevel{junior, confirmé}]\label{iq:iv2:rates-and-fx:17}
Trois obligations au pair à deux, cinq et dix ans ont un DV01, pour 100~\$, de 1.886, 4.452 et 8.111 cents. Vous vendez à découvert 100~\$ de l'obligation à cinq ans et achetez celles à deux et dix ans, de sorte que le trade n'ait ni flux de trésorerie net ni DV01. Quelle est la taille de chaque jambe~? Quel est votre profit si le rendement à cinq ans monte de 10 points de base et que les deux autres ne bougent pas~? Sur quoi pariez-vous~?
\end{question}

\begin{question}[$\star\star$ \iqroles{QT, Trader, Risque} \iqfirm{banque, hedge fund} \iqlevel{junior}]\label{iq:iv2:rates-and-fx:18}
L'obligation à deux ans a un DV01 de 190~\$ pour 1 million de dollars, et celle à dix ans de 800~\$ pour 1 million. Montez un steepener neutre en DV01 avec 10 millions de dollars de l'obligation à dix ans. Que gagnez-vous si le taux à deux ans baisse de 5 points de base et celui à dix ans monte de 5~? Et si les deux montent de 10~?
\end{question}

\begin{question}[$\star\star$ \iqroles{QT, Trader, QR, Risque} \iqfirm{banque, hedge fund} \iqlevel{junior, confirmé}]\label{iq:iv2:rates-and-fx:19}
Les variations de rendement ont une volatilité annualisée de 80 points de base pour l'obligation à deux ans et de 75 pour celle à dix ans, avec une corrélation de 0.7. Vous détenez 10 millions de dollars de l'obligation à dix ans, de DV01 8\,000~\$, et vous couvrez avec celle à deux ans. Quel DV01 de couverture donne une régression~? Comment le risque résiduel se compare-t-il à une couverture neutre en DV01~?
\end{question}

\begin{question}[$\star\star$ \iqroles{QR, QT, Risque} \iqfirm{banque, hedge fund} \iqlevel{junior}]\label{iq:iv2:rates-and-fx:20}
Vous recevez le taux fixe de 4~\% sur un swap dont il reste deux paiements annuels, sur un notionnel de 100~\$. La courbe est désormais plate à 3~\%. Quelle est la valeur du swap~? Calculez-la de deux façons~: comme une obligation moins le pair, puis comme l'écart de taux multiplié par l'annuité.
\end{question}

\begin{question}[$\star\star$ \iqroles{QR, QT, Trader} \iqfirm{banque, hedge fund} \iqlevel{junior, confirmé}]\label{iq:iv2:rates-and-fx:21}
Un contrat à terme sur taux à trois mois expire dans cinq ans, et il fixe un taux de l'année 5 à l'année 5.25. La volatilité normale de ce taux est de 1~\% par an. Quel est l'ajustement de convexité entre le taux du contrat à terme et le taux forward, selon la règle habituelle~? Lequel est le plus élevé, et pourquoi~?
\end{question}

\begin{question}[$\star\star$ \iqroles{QT, Trader, QR} \iqfirm{banque, hedge fund} \iqlevel{junior, confirmé}]\label{iq:iv2:rates-and-fx:22}
Un contrat à terme sur obligation se négocie à 100. Le vendeur peut livrer l'une de trois obligations, de facteurs de concordance 0.92, 1.10 et 1.25, et de prix 93.0, 110.8 et 125.6. Laquelle est la moins chère à livrer~? Pourquoi la réponse change-t-elle quand les rendements bougent~? Ignorez le carry.
\end{question}

\begin{question}[$\star\star$ \iqroles{QT, Trader, QR} \iqfirm{banque, hedge fund} \iqlevel{junior, confirmé}]\label{iq:iv2:rates-and-fx:23}
Vous achetez une obligation au pair à cinq ans à un rendement de 4~\%, financée au taux de repo de 3~\%. Le rendement à quatre ans est de 3.8~\%, et une obligation à quatre ans a une duration de 3.6. Quel est votre carry, et votre roll-down sur un an, si la courbe ne bouge pas~? De combien le rendement à quatre ans peut-il monter à l'horizon avant que vous ne perdiez de l'argent~?
\end{question}

\begin{question}[$\star\star$ \iqroles{QR, QD, Risque} \iqfirm{banque, hedge fund} \iqlevel{junior, confirmé}]\label{iq:iv2:rates-and-fx:24}
Dans le modèle de Vasicek, le taux court part de 3~\%, revient vers 5~\% avec une vitesse de 0.5, et a une volatilité de 1~\%. Quelles sont la moyenne et l'écart-type du taux court dans deux ans~? Quelle est la probabilité d'un taux négatif à cette date~? Cela rend-il le modèle inutilisable~?
\end{question}

\begin{question}[$\star\star$ \iqroles{QR, QD} \iqfirm{banque, hedge fund} \iqlevel{junior, confirmé}]\label{iq:iv2:rates-and-fx:25}
Un taux court sans dérive reste plat en espérance à un niveau $f$ (Ho-Lee à dérive nulle), et il a une volatilité normale de 1~\%. Quel est le rendement zéro-coupon à dix ans, en points de base sous $f$~? Dérivez le $\sigma^2T^2/6$. Pourquoi un taux volatil fait-il baisser les rendements longs~?
\end{question}

\section{Options sur les taux et le change}

\begin{question}[$\star\star$ \iqroles{QR, QT, Trader} \iqfirm{banque, hedge fund} \iqlevel{junior}]\label{iq:iv2:rates-and-fx:26}
Pricez un caplet sur un taux à trois mois fixé dans un an, avec un taux forward de 4~\%, un strike de 4~\%, une volatilité lognormale de 20~\% et un facteur d'actualisation de 0.95 jusqu'à la date de paiement, par unité de notionnel. Quelle est cette même volatilité en points de base, et quel est le prix du caplet selon Bachelier~? Quelle est la volatilité normale si le taux vaut 2~\%, à la même volatilité lognormale~?
\end{question}

\begin{question}[$\star\star$ \iqroles{QR, QT, Trader} \iqfirm{banque, hedge fund} \iqlevel{junior}]\label{iq:iv2:rates-and-fx:27}
Une swaption à la monnaie à un an porte sur un swap dont l'annuité vaut 4.5. La volatilité normale est de 80 points de base. Quelle est sa valeur, en pourcentage du notionnel, selon la règle empirique~? Pourquoi l'annuité la multiplie-t-elle~? De combien le prix varie-t-il pour un point de base de volatilité~?
\end{question}

\begin{question}[$\star\star$ \iqroles{QR, QT, Trader} \iqfirm{banque, hedge fund} \iqlevel{junior}]\label{iq:iv2:rates-and-fx:28}
Une swaption payeuse et une swaption receveuse à un an ont le même strike de 5~\%. Le taux de swap forward est de 4~\%, et l'annuité vaut 4.5. Quelle est la valeur de la payeuse moins la receveuse~? Démontrez-le. Quelle est cette position, et quel modèle avez-vous utilisé~?
\end{question}

\begin{question}[$\star\star$ \iqroles{QT, Trader, QR} \iqfirm{banque, hedge fund} \iqlevel{junior, confirmé}]\label{iq:iv2:rates-and-fx:29}
Le spot EURUSD vaut 1.10, le taux en euros à un an est de 3~\%, et le forward à un an se négocie à 1.1245. Quel taux en dollars le marché des swaps implique-t-il, en taux simples~? Le taux monétaire en dollars est de 5~\%. Qu'est-ce que cet écart vous dit-il~? Pourquoi l'arbitrage ne le fait-il pas disparaître~?
\end{question}

\begin{question}[$\star\star$ \iqroles{QR, QT, Trader} \iqfirm{banque, hedge fund} \iqlevel{junior}]\label{iq:iv2:rates-and-fx:30}
EURUSD vaut 1.10. Un call sur euro à un an, à la monnaie spot, a un strike de 1.10, une volatilité de 10~\%, un taux en dollars de 5~\% et un taux en euros de 3~\%, à capitalisation continue. Quel est le forward~? Quelles sont la valeur du call et celle du put, en dollars par euro~? Quels sont le delta spot et le delta forward~?
\end{question}

\begin{question}[$\star\star$ \iqroles{QR, QT, Trader} \iqfirm{banque, hedge fund} \iqlevel{junior, confirmé}]\label{iq:iv2:rates-and-fx:31}
Un call sur euro à un an, avec EURUSD à 1.10, un strike de 1.10, une volatilité de 10~\%, un taux en dollars de 5~\% et un taux en euros de 3~\%, à capitalisation continue, vaut 0.0536~\$ par euro. Quelle est cette valeur en euros, et en pourcentage du notionnel en euros et du notionnel en dollars~? Montrez que la même option évaluée côté dollar, comme un put en dollars sur le taux de change en euros par dollar à un strike de 0.909, donne la même valeur. Quelle doit être la volatilité, et pourquoi~?
\end{question}

\begin{question}[$\star\star$ \iqroles{QT, Trader, QR} \iqfirm{banque, hedge fund} \iqlevel{junior, confirmé}]\label{iq:iv2:rates-and-fx:32}
Un desk de change cote le smile à un an avec une volatilité à la monnaie de 10~\%, un risk reversal à 25 delta de $-2~\%$ et un butterfly à 25 delta de 0.5~\%. Quelles sont les volatilités du call et du put à 25 delta~? Avec un forward de 1.1222, un taux en euros de 3~\%, et en utilisant le delta spot, à quels strikes sont-elles associées~?
\end{question}

\begin{question}[$\star\star$ \iqroles{QT, Trader, QR, Risque} \iqfirm{banque, hedge fund} \iqlevel{junior, confirmé}]\label{iq:iv2:rates-and-fx:33}
EURUSD a une volatilité de 10~\%, et USDJPY une volatilité de 9~\%, avec une corrélation de 0.5 entre leurs variations logarithmiques. Quelle est la volatilité de EURJPY~? Et si la corrélation vaut $-0.5$~? Comment utilisez-vous ce résultat à l'envers, et que signifie-t-il pour le trading de trois devises~?
\end{question}

\section{Modèles, produits et conventions}

\begin{question}[$\star\star\star$ \iqroles{QR, QT, Trader} \iqfirm{banque, hedge fund} \iqlevel{junior, confirmé}]\label{iq:iv2:rates-and-fx:34}
La jambe variable d'un swap overnight compose un taux quotidien sur une semaine~: 4~\%, 4~\%, 5~\% et 4~\% du lundi au jeudi, puis 4~\% le vendredi pour trois jours, en base actual/360. Quel taux la jambe variable verse-t-elle pour la semaine~? Comment se compare-t-il à une moyenne simple~? Contrairement à un taux term fixé à l'avance, quand est-il connu~? Quel effet cela a-t-il sur un caplet porté par cette jambe~?
\end{question}

\begin{question}[$\star\star\star$ \iqroles{QR, QD} \iqfirm{banque, hedge fund} \iqlevel{junior, confirmé}]\label{iq:iv2:rates-and-fx:35}
Écrivez le prix du zéro-coupon de Vasicek sous forme fermée, avec une vitesse de 0.5, une moyenne de 5~\%, une volatilité de 1~\% et un taux initial de 3~\%. Quel est le zéro à cinq ans, et quel est son rendement~? Dérivez $B(T)$. Vérifiez le prix par une simulation de Monte-Carlo de l'actualisation. Comment le rendement se comporte-t-il à très longue échéance~?
\end{question}

\begin{question}[$\star\star\star$ \iqroles{QR, QD} \iqfirm{banque, hedge fund} \iqlevel{confirmé}]\label{iq:iv2:rates-and-fx:36}
Dans le modèle à un facteur de Hull--White, avec une vitesse de 0.1 et une volatilité de 1~\%, la courbe de marché est plate à 4~\%, à capitalisation continue. Trouvez la dérive $\theta(t)$ qui l'ajuste. Pourquoi la moyenne constante de Vasicek ne suffit-elle pas~? Vérifiez par simulation que le modèle restitue le zéro à dix ans. Que vaut $f(t,T)$ par la suite~?
\end{question}

\begin{question}[$\star\star\star$ \iqroles{QR, QT, Trader} \iqfirm{banque, hedge fund} \iqlevel{confirmé}]\label{iq:iv2:rates-and-fx:37}
Deux taux forward annuels valent initialement 4~\%, avec des volatilités lognormales de 20~\%, et ils sont fixés dans un an, avec des facteurs d'actualisation de 0.96 et 0.92. Prenez le taux de swap comme leur moyenne. Un cap comporte deux caplets de strike 4~\%, et une swaption payeuse a le même strike. Quelle est la valeur de chacun, en points de base, avec des corrélations de 1, 0.5 et 0~? Pourquoi la swaption ne vaut-elle jamais plus que le cap~? Quel marché cela relie-t-il~?
\end{question}

\begin{question}[$\star\star\star$ \iqroles{QR, QT, Risque} \iqfirm{banque, hedge fund} \iqlevel{junior, confirmé}]\label{iq:iv2:rates-and-fx:38}
Une obligation à cinq ans, à coupon de 5~\%, peut être rachetée au pair par l'émetteur dans un an. Prenez un rendement plat, une volatilité normale du rendement de 100 points de base, et évaluez le call comme une option, dans un an, sur l'obligation à quatre ans. Quelles sont les valeurs de l'obligation bullet, du call et de l'obligation rachetable à des rendements de 4~\%, 5~\% et 6~\%~? Quelles sont leurs durations~? Qu'est-ce que la convexité négative, et qui est short de cette convexité~?
\end{question}

\begin{question}[$\star\star\star$ \iqroles{QT, Trader, QR} \iqfirm{banque, hedge fund} \iqlevel{confirmé}]\label{iq:iv2:rates-and-fx:39}
Un contrat à terme peut livrer une obligation à cinq ans ou à vingt ans, à coupon de 6~\%, chacune avec un facteur de concordance de 1, puisque les facteurs sont calculés à un rendement de 6~\%. Le rendement est à 6~\%, la courbe est plate, et à la livraison il est normal, avec un écart-type de 100 points de base. Quel est le prix juste du contrat à terme, comme moyenne du moins cher des deux~? De combien le choix du short retranche-t-il au prix~? Quelle obligation est la moins chère à livrer si les rendements montent~?
\end{question}

\begin{question}[$\star\star\star$ \iqroles{QR, QT, Risque} \iqfirm{banque, hedge fund} \iqlevel{junior, confirmé}]\label{iq:iv2:rates-and-fx:40}
Les variations de rendement à 2, 5, 10 et 30 ans ont des volatilités de 80, 85, 75 et 65 points de base. Leurs corrélations valent 0.85, 0.7 et 0.6 pour le 2 ans avec les autres, 0.9 et 0.8 pour le 5 ans avec les deux plus longs, et 0.9 pour le 10 ans avec le 30 ans. Quelle part de la variance expliquent les trois premières composantes principales~? À quoi ressemblent-elles~? Comment les utiliseriez-vous pour couvrir le risque~?
\end{question}

\begin{question}[$\star\star\star$ \iqroles{QR, QT} \iqfirm{banque, hedge fund} \iqlevel{confirmé}]\label{iq:iv2:rates-and-fx:41}
Un contrat verse un taux annuel qui se fixe dans cinq ans, et il paie à la date du fixing, et non un an plus tard. Le forward vaut 4~\%, avec une volatilité lognormale de 20~\%. Quel est l'ajustement de convexité~? Dérivez-le à partir d'un changement de probabilité, et comparez la valeur exacte à la règle du premier ordre. Pourquoi apparaît-il~?
\end{question}

\begin{question}[$\star\star\star$ \iqroles{QR, QT} \iqfirm{banque, hedge fund} \iqlevel{confirmé}]\label{iq:iv2:rates-and-fx:42}
Un swap CMS verse le taux de swap à dix ans, fixé dans cinq ans, en un seul paiement à la date du fixing. Le taux de swap forward vaut 4~\%, avec une volatilité lognormale de 20~\%. Utilisez une annuité à rendement plat $A(y)$ pour estimer l'ajustement. Que donne la règle du premier ordre~? Que donne une simulation sous la probabilité annuité~? Dans quel sens joue l'ajustement, et pourquoi dépend-il de tout le smile~?
\end{question}

\begin{question}[$\star\star\star$ \iqroles{QR, QT} \iqfirm{banque, hedge fund} \iqlevel{confirmé}]\label{iq:iv2:rates-and-fx:43}
Un contrat payé en dollars porte sur un indice yen à 100 et verse un dollar par point d'indice. L'indice a une volatilité de 20~\%, la valeur en dollars d'un yen a une volatilité de 10~\%, et leur corrélation vaut $-0.3$. Les taux sont nuls. Quel est le forward quanto~? Dérivez l'ajustement de dérive, et vérifiez-le par simulation. Pourquoi le signe va-t-il dans ce sens~?
\end{question}

\begin{question}[$\star\star\star$ \iqroles{QT, Trader, QR} \iqfirm{banque, hedge fund} \iqlevel{confirmé}]\label{iq:iv2:rates-and-fx:44}
Une option de change à un an sur un forward de 1.1222 a une volatilité de 10~\%. Quel est le strike auquel un straddle a un delta nul, lorsque le delta n'intègre pas la prime~? Et lorsqu'il l'intègre~? Pourquoi certaines paires de devises utilisent-elles la forme ajustée de la prime~? Qu'est-ce qui va mal si vous confondez les deux~?
\end{question}

\begin{question}[$\star\star\star$ \iqroles{QD, QR, Risque} \iqfirm{banque, hedge fund} \iqlevel{junior, confirmé}]\label{iq:iv2:rates-and-fx:45}
Les taux zéro-coupon à capitalisation continue à 1, 2, 3 et 5 ans valent 3~\%, 4~\%, 3.8~\% et 4.2~\%. Vous interpolez linéairement le taux zéro. Quel est le forward instantané de part et d'autre de deux ans~? Et avant et après trois ans~? Que donne l'interpolation log-linéaire des facteurs d'actualisation~? Laquelle utiliseriez-vous pour la couverture, et pourquoi~?
\end{question}

\begin{question}[$\star\star\star$ \iqroles{QR, QT, Trader} \iqfirm{banque, hedge fund} \iqlevel{junior, confirmé}]\label{iq:iv2:rates-and-fx:46}
Une obligation à cinq ans au pair a une duration de 4.45. Avec une volatilité normale du rendement de 80 points de base, quelle est la volatilité du prix de l'obligation~? Évaluez un call à un an à la monnaie forward sur celle-ci selon la règle empirique. Comparez ce résultat à une swaption receveuse à un an, à 80 points de base, sur une annuité de 4.5. Pourquoi sont-elles presque égales~? Laquelle des deux porte le risque supplémentaire~?
\end{question}

\begin{question}[$\star\star\star$ \iqroles{QT, Trader, QR} \iqfirm{banque, hedge fund} \iqlevel{junior, confirmé}]\label{iq:iv2:rates-and-fx:47}
Une obligation cotée 100 verse un coupon de 2 dans le prochain semestre, et elle est la moins chère à livrer dans un contrat à terme~; le prix du contrat à terme multiplié par le facteur de concordance vaut 99.5 à la livraison, dans un semestre. Quel est le taux de repo implicite, en base simple~? Le taux de repo de marché pour cette obligation est de 2.5~\%. Quel trade mettez-vous en place, et qu'est-ce qui peut le faire échouer~?
\end{question}

\section{Trois à construire}

\begin{question}[$\star\star\star\star$ \iqroles{QR, QD} \iqfirm{banque, hedge fund} \iqlevel{confirmé}]\label{iq:iv2:rates-and-fx:48}
Dans le modèle de Hull-White, avec une vitesse de 0.1, une volatilité de 1~\% et une courbe plate à 4~\% à capitalisation continue, pricez une swaption payeuse à la monnaie à un an sur un swap annuel à cinq ans. Utilisez la décomposition de Jamshidian, et vérifiez par simulation. Quel est le strike~? Pourquoi la décomposition fonctionne-t-elle à un facteur et échoue-t-elle à deux~?
\end{question}

\begin{question}[$\star\star\star\star$ \iqroles{QR, QT, QD} \iqfirm{banque, hedge fund} \iqlevel{confirmé}]\label{iq:iv2:rates-and-fx:49}
Un smile de change à un an présente les cotations suivantes~: une volatilité à la monnaie de 10~\%, au strike de 1.1278~; un put à 25 delta à 11.5~\%, au strike de 1.0482~; et un call à 25 delta à 9.5~\%, au strike de 1.1992. Le forward vaut 1.1222, et le taux en dollars est de 5~\%. Dérivez les poids vanna-volga permettant de répliquer une nouvelle option à partir de ces trois cotations. Que donnent-ils pour des strikes de 1.10 et de 1.25~? Pourquoi les piliers sont-ils retrouvés exactement~? Où cesseriez-vous de lui faire confiance~?
\end{question}

\begin{question}[$\star\star\star\star$ \iqroles{QR, QD} \iqfirm{banque, hedge fund} \iqlevel{confirmé}]\label{iq:iv2:rates-and-fx:50}
Dans le modèle de Hull--White, avec une vitesse de 0.1, une volatilité de 1~\% et une courbe plate à 4~\% à capitalisation continue, une swaption bermudéenne payeuse peut être exercée aux années 1, 2, 3 et 4, dans un swap à paiements annuels qui se termine à l'année 5, à un strike égal au taux juste du swap commençant à l'année 1. Remontez-la sur une grille. Quelles sont les quatre valeurs européennes et la valeur bermudéenne~? Quelle est la valeur du switch~? Que fait la vitesse de retour à la moyenne à cette valeur du switch, et pourquoi~?
\end{question}
